% =====================================================================
\section{Conclusion}
\label{sec:conclusion}
% =====================================================================

We have presented a content-addressing system for formal mathematics,
implemented as a Lean~4 library.  By hashing canonicalized type
expressions, we assign declarations stable identifiers determined by
their mathematical content rather than their human-chosen names.

Content-addressing does not require full normalization to be useful.
Even partial canonicalization (Level~0) provides meaningful stability
and cross-project discoverability.  The decentralized registry
architecture, made possible by the monotonicity of mathematical proof,
requires no infrastructure beyond Git---though significant engineering
work remains to make it practical at ecosystem scale.

The system is a working prototype, and important challenges
remain---most fundamentally, the gap between our canonicalization levels
and full definitional equality.  We hope this work demonstrates that
content-addressing is a productive lens through which to view the
organization of formal mathematical knowledge, and that the tools
presented here are a useful step toward a more interconnected ecosystem
for formalized mathematics.

\simone{Separate language-agnostic parts (registry format, discovery protocol, Merkle DAG) from Lean-specific parts (canonicalization, \texttt{Expr} serialization, \texttt{whnf} batching). Former could become standard for other proof assistants.}


% =====================================================================
% References
% =====================================================================

